\section{Method}
\label{sec:method}

Every symbol, with its dimensions and their meaning, is also listed in Table~\ref{tab:notation} (Appendix~\ref{app:notation}).

\begin{figure*}[t]
\centering
\resizebox{\textwidth}{!}{%
\begin{tikzpicture}[
  font=\small,
  box/.style={draw, rounded corners=2pt, minimum height=0.9cm, align=center, inner sep=4pt},
  fz/.style={box, fill=frozen!15, draw=frozen, thick},
  tr/.style={box, fill=trained!15, draw=trained, thick},
  tok/.style={draw, minimum width=0.38cm, minimum height=0.42cm, inner sep=0pt},
  arr/.style={-{Stealth[length=2.2mm]}, thick}]
  % inputs and the two streams
  \node[box, fill=gray!10, minimum width=1.9cm] (x) at (0,0) {Meeting\\audio $\bm{x}$};
  \node[fz, minimum width=2.7cm] (enc) at (3.2,1.0) {Audio encoder\\(\S\ref{sec:method-task})};
  \node[fz, minimum width=2.7cm] (dia) at (3.2,-1.0) {Streaming diarizer\\$S{=}8$ slots};
  \node[tr, minimum width=2.3cm] (proj) at (6.5,-1.0) {Speaker projector\\(App.~\ref{app:model})};
  \draw[arr] (x.east) -- ++(0.3,0) |- (enc.west);
  \draw[arr] (x.east) -- ++(0.3,0) |- (dia.west);
  \draw[arr] (dia) -- node[above]{$\bm{P}$} (proj);
  % interleaved sequence
  \foreach \i/\c in {0/frozen,1/frozen,2/frozen,3/frozen,4/frozen,5/spkc,6/frozen,7/frozen,8/frozen,9/frozen,10/frozen,11/spkc} {
    \node[tok, fill=\c!35] (t\i) at (8.9+\i*0.4,0) {};
  }
  \node (dots) at (8.9+12*0.4+0.05,0) {$\cdots$};
  \draw[decorate, decoration={brace, amplitude=4pt}] (t0.north west) -- node[above=4pt, font=\small]{$\bm{A}_1$: $K{=}5$ acoustic tokens} (t4.north east);
  \node[below=2pt of t5, font=\small] {$\bm{s}_1$};
  \node[below=2pt of t11, font=\small] {$\bm{s}_2$};
  \node[font=\small] at (11.2,-1.25) {interleaved input $\bm{E}$ (\S\ref{sec:method-task})};
  \draw[arr] (enc.east) -| node[pos=0.3, above, font=\small]{acoustic tokens $\bm{a}_i$} (t8.north);
  \draw[arr] (proj.east) -- node[above, font=\small]{$\bm{s}_j$} ++(0.75,0) |- (t0.west);
  % LLM
  \node[tr, minimum height=1.9cm, minimum width=2.0cm] (llm) at (15.6,0) {Thinker LLM\\+ LoRA\\(\S\ref{sec:method-out})};
  \draw[arr] (dots.east) -- (llm.west);
  % output segments (second row)
  \node[box, fill=white, anchor=west] (o1) at (6.2,-2.6)
    {$\spk{2}\,\tim{5}$\ \zh{我想知道会议的时间}\ $\tim{48}$};
  \node[box, fill=white, anchor=west] (o2) at (6.2,-3.7)
    {$\spk{3}\,\tim{42}$\ \zh{对，没错。}\ $\tim{56}$};
  \draw[arr] (llm.south) -- ++(0,-0.85) -| ([xshift=-0.3cm]o1.west) -- (o1.west);
  \draw[arr] ([xshift=-0.3cm]o1.west) |- (o2.west);
  \node[font=\small, anchor=west] at (6.2,-4.45) {who, when, what (Eq.~\ref{eq:serial})};
  % talker
  \node[tr, minimum width=2.4cm] (talk) at (14.9,-3.15) {Talker, voice $\bm{v}_k$\\playback alignment\\(\S\ref{sec:method-plem})};
  \draw[arr] (o1.east) -- ++(0.25,0) |- (talk.west);
  \draw[arr] (o2.east -| o1.east) ++(0,0) -- ++(0.25,0) |- (talk.west);
  \draw[thick] (o2.east) -- (o2.east -| o1.east);
  \node[box, fill=gray!10, minimum width=1.7cm] (y) at (17.6,-3.15) {Translated\\speech};
  \draw[arr] (talk) -- (y);
  % voice prompts
  \draw[arr, dashed, spkc!80!black] (dia.south) -- ++(0,-3.75) -| node[pos=0.25, below, font=\small]{voice prompt $\bm{v}_k$: solo regions of slot $k$ (Eq.~\ref{eq:alone})} (talk.south);
  % streaming loop
  \draw[arr, dotted, trained!80!black] (llm.north) -- ++(0,0.5) -- node[above, font=\small]{streaming: read every $\Delta$\,s, write or $\wait$ on deadlines (\S\ref{sec:method-simul})} ++(-7.4,0);
\end{tikzpicture}}
\caption{\sys{}. A frozen audio encoder turns the meeting into acoustic tokens (what is said); a frozen streaming
diarizer gives slot activities $\bm{P}$ that a trained projector turns into speaker tokens $\bm{s}_j$ (who says it).
One speaker token follows every block of $K{=}5$ acoustic tokens. The Thinker writes speaker-attributed, time-stamped
translation segments, offline or read by read, and a Talker speaks each segment in the voice of its speaker slot;
its playback is trained by playback alignment (\S\ref{sec:method-plem}) with the Thinker fixed. Blue: frozen; orange:
trained; dotted: the streaming loop.}
\label{fig:arch}
\end{figure*}


\subsection{Task and input: who said what}
\label{sec:method-task}
\label{sec:method-audio}
\label{sec:method-spk}
\label{sec:method-fuse}

\paragraph{Task.} The input is a single-channel recording $\bm{x}\in\R^{16000T}$ of $T$ seconds (16\,kHz samples) in
which an unknown number of people talk, sometimes at once. The output is a list of $N_\mathcal{Y}$ segments
$\mathcal{Y}=\big((k_n,a_n,b_n,\bm{y}^{(n)})\big)_{n=1}^{N_\mathcal{Y}}$: $k_n\in\{1,\dots,S\}$ is the speaker label of
segment $n$ (arbitrary, but one person keeps one label for the whole meeting), $a_n<b_n\in[0,T]$ are the start and end
of its source speech in seconds, and $\bm{y}^{(n)}$ is its translation, a sequence of $|\bm{y}^{(n)}|$ text tokens; in
speech output $\bm{y}^{(n)}$ is spoken in the voice of speaker $k_n$. In streaming mode segments are produced while
$\bm{x}$ is recorded.

\paragraph{Two streams.} The frozen audio encoder of Qwen3-Omni, whose LLM (the ``Thinker'') already translates
single-speaker speech~\citep{xu2025qwen3omni}, maps the recording to acoustic tokens
$\bm{A}=[\bm{a}_1;\dots;\bm{a}_N]\in\R^{N\times d}$: row $i$ is the token for the 80\,ms of audio centred at $\bar{t}_i$
(so $N\approx T/0.08$), and its $d=2048$ columns are the Thinker's hidden dimensions. The tokens carry no memory of who
spoke earlier. A frozen Streaming Sortformer~\citep{medennikov2025streaming} keeps every voice in the slot it got when
it first appeared and returns $\bm{P}\in[0,1]^{F\times S}$: row $f$ is a frame of 80\,ms at time $t_f$, column $k$ is
one of $S=8$ speaker slots, and $P_{f,k}$ is the probability that slot $k$ is speaking in frame $f$; overlapped speech
gives several large entries in one row. Every $K=5$ rows of $\bm{P}$, resampled to the token centres, become one
speaker token $\bm{s}_j\in\R^{d}$ (one per 0.4\,s, in the Thinker's hidden space) through a strided convolution and a
two-layer perceptron whose output layer starts at zero (Appendix~\ref{app:model}); a streaming frame is never revised
(Appendix~\ref{app:causal}).

\paragraph{Interleaved fusion.} Following G-STAR~\citep{peng2026gstar}, each speaker token follows the $K$ acoustic
tokens it describes: $\bm{E}=[\bm{A}_1;\bm{s}_1;\dots;\bm{A}_M;\bm{s}_M]\in\R^{(N+M)\times d}$, where
$\bm{A}_j\in\R^{K\times d}$ is the $j$-th block of acoustic tokens and $M=\lceil N/K\rceil$ the number of blocks.
$\bm{E}$ fills the audio placeholder of the chat template after a text instruction $\bm{\iota}$ (Appendix~\ref{app:prompt}).

\subsection{Speaker-attributed, time-stamped output}
\label{sec:method-out}
\label{sec:method-train}

\paragraph{Serialisation.} We add 251 time tokens $\tim{0},\dots,\tim{250}$ ($\tim{m}$ marks $0.08m$\,s, up to
20\,s) and 8 speaker tokens $\spk{1},\dots,\spk{8}$ (one per slot) to the vocabulary, with trainable input and output
rows, and write a window's segments speaker first, ordered by start as in serialized output training~\citep{kanda2020sot}:
\begin{equation}
\bm{y}=\mathop{\Vert}_{n=1}^{N_\mathcal{Y}}\big[\spk{k_n}\,\tim{\alpha_n}\,\bm{y}^{(n)}\,\tim{\beta_n}\big],
\label{eq:serial}
\end{equation}
where $\Vert$ concatenates, and $\alpha_n,\beta_n\in\{0,\dots,250\}$ are $a_n,b_n$ in 80\,ms steps from the window's
start. Targets use the diarizer's slots: per meeting, each reference speaker is mapped to the slot that overlaps most
with that speaker's speech over the whole meeting (Hungarian matching;~\citealp{kuhn1955hungarian}), so the model
learns to copy the slot rather than invent a label per window (Appendix~\ref{app:model}).

\paragraph{Training.} The Thinker with parameters $\theta$ defines
$p_\theta(\bm{y}\mid\bm{E},\bm{\iota})=\prod_{u=1}^{|\bm{y}|}p_\theta(y_u\mid\bm{y}_{<u},\bm{E},\bm{\iota})$ over the
$|\bm{y}|$ tokens of Eq.~\eqref{eq:serial}, and is trained with the weighted cross-entropy
$\mathcal{L}=-\Omega^{-1}\sum_u\omega(y_u)\log p_\theta(y_u\mid\bm{y}_{<u},\bm{E},\bm{\iota})$, where the token weight
$\omega(y_u)$ is 2 for speaker tokens, 1.5 for time tokens and 1 otherwise (one wrong speaker or time token can
misattribute a sentence) and $\Omega=\sum_u\omega(y_u)$ normalises. Trained are the projector, the new rows and LoRA
adapters~\citep{hu2022lora} on all 48 layers (Table~\ref{tab:hyper}); \sys{} (offline) decodes a window at once.

\subsection{Streaming interpretation}
\label{sec:method-simul}

\paragraph{Reading and writing.} The model reads every $\Delta=1$\,s, at times $t_r=\min(r\Delta,T)$ for reads
$r=1,\dots,R$, $R=\lceil T/\Delta\rceil$. At read $r$ it sees the audio $\bm{x}_{\le t_r}$ and a causal speaker stream:
the low-latency diarizer finalises a frame at most $\gamma=1.04$\,s later, so frames after $t_r-\gamma$ are zeroed,
$\bm{P}^{(r)}\in[0,1]^{F\times S}$ with $P^{(r)}_{f,k}=P_{f,k}\mathds{1}[t_f\le t_r-\gamma]$, and $\bm{E}^{(r)}$ is built
from $\bm{x}_{\le t_r}$ and $\bm{P}^{(r)}$ as above. With $\hat{\bm{y}}^{(<r)}$ the tokens committed before read $r$, the
model continues greedily from $(\hat{\bm{y}}^{(<r)},\bm{E}^{(r)},\bm{\iota})$ until a stop token $\eta_r$, the wait
token $\wait$ for $r<R$ and $\eos$ for $r=R$; what precedes it is committed and never revised. Scoring times commits by
the actual release time of each frame (Appendix~\ref{app:chunk}).

\paragraph{Deadline targets.} Each reference translation is split at punctuation into chunks, and each chunk is timed by
the source it covers (word alignments, or the same fraction of the utterance where none exist; Appendix~\ref{app:model}).
A chunk $j$ whose source ends at $e_j$ seconds becomes due one interpreting lag $\epsilon=1$\,s later, at the read
$r(j)=\min(\{r: t_r\ge e_j+\epsilon\}\cup\{R\})$. The target of read $r$ is the chunks due at $r$, serialised as in
Eq.~\eqref{eq:serial}, then $\eta_r$; each read is one training row whose prefix is the targets of all earlier reads. A
deadline model is trained on these rows with the same loss, starting from \sys{} (offline). One such model, trained on
both directions (\S\ref{sec:setup}), is the Thinker of \sys{}; with the base Talker it is the \emph{deadline baseline}
(B). \sys{}-Simul, used in the overlap study, has one deadline model per direction.

\subsection{Speaking in each speaker's voice}
\label{sec:method-voice}

\paragraph{Voices.} Runs of at least 0.4\,s in which slot $k$ alone is active (Eq.~\eqref{eq:alone},
Appendix~\ref{app:model}) are concatenated into a voice prompt $\bm{v}_k$, at most 8\,s of audio, in streaming mode only
from audio heard before the slot's first commitment. The UniSS Talker~\citep{uniss2025} encodes $\bm{v}_k$ into 32
BiCodec global tokens~\citep{wang2025spark} and speaks each of the slot's text blocks as a sequence of BiCodec semantic
tokens, one per 20\,ms of speech. Each speaker has an output stream of its own, in which a block starts at its release
time or when the previous block ends (Eq.~\eqref{eq:layout}). The base Talker plays every committed block in full at a
fixed rate, so a stream that receives more speech than its source builds a queue. The playback layer instead gives the
Talker a base rate that rises with the backlog: the source's pace plus three rate steps per second of backlog, at most
six (Appendix~\ref{app:plem}).

\subsection{Playback alignment: a CRF over when each block is heard}
\label{sec:method-plem}

\paragraph{Queue and delay.} Consider one speaker's output stream with $I$ blocks, $i=1,\dots,I$. Block $i$ holds
$W_i\in\mathbb{N}$ target units (words for English, characters for Chinese) and is released to the Talker at time
$\nu_i$ in seconds (\S\ref{sec:method-online}). The base Talker proposes a finite set $C_i$ of $|C_i|=8$ renderings of
it (samples at two speaking rates); a candidate $c\in C_i$ is a sequence $\bm{a}(c)$ of semantic tokens and lasts
$D(c)=0.02\,|\bm{a}(c)|$ seconds. Let the \emph{backlog} $h_i$ be how many seconds block $i$ waits after $\nu_i$ for
earlier speech of its speaker. The first block does not wait, $h_1=0$, and block $i$ ends at $\nu_i+h_i+D(c_i)$, so the
next block waits for whatever of that remains at its own release:
\begin{equation}
h_{i+1}=g_i(h_i,c_i)\coloneqq\Pi_{\mathcal{H}}\big(\nu_i+h_i+D(c_i)-\nu_{i+1}\big),\!
\label{eq:backlog}
\end{equation}
where $\Pi_{\mathcal{H}}$ rounds $\max(0,\cdot)$ up to the grid $\mathcal{H}=\{0,\delta,\dots,\bar{H}\}$ of
$|\mathcal{H}|=101$ backlog values ($\delta=80$\,ms, cap $\bar{H}=8$\,s), so that
$g_i:\mathcal{H}\times C_i\to\mathcal{H}$. Block $i$ starts playing $h_i$ after its release; spreading its $W_i$ units
evenly over its $D(c_i)$ seconds of playback, the $k$-th unit has been heard $h_i+\frac{k}{W_i}D(c_i)$ after $\nu_i$.
Summed over the block, the delay at which the listener hears its units is, in unit-seconds,
\begin{equation}
L_i(h_i,c_i)=W_i\,h_i+\tfrac{W_i+1}{2}\,D(c_i).
\label{eq:lat}
\end{equation}
A long rendering therefore costs twice: its own units are spread over more time, and through Eq.~\eqref{eq:backlog}
every later block waits longer.

\paragraph{From scores to a CRF.} A \emph{playback plan} is a sequence $\bm{z}=(z_1,\dots,z_I)$ of states
$z_i=(h_i,c_i)\in\mathcal{H}\times C_i$, one per block (a backlog value and a rendering); it is feasible,
$\bm{z}\in\mathcal{Z}(\bm{x})$, if $h_1=0$ and every backlog follows Eq.~\eqref{eq:backlog}, where $\bm{x}$ now denotes
the stream (block texts, release times $\nu_{1:I}$, candidate sets $C_{1:I}$). Each state is scored by a real-valued
local potential $\phi^\lambda_i:\mathcal{H}\times C_i\to\R$,
\begin{equation}
\phi^\lambda_i(h,c)=\underbrace{0.5\,\bar{\ell}(c)+5f(c)}_{\text{quality}}-\underbrace{\lambda L_i(h,c)}_{\text{delay}},
\label{eq:plpot}
\end{equation}
where $\bar{\ell}(c)\le0$ is the Talker's mean log-probability per token of $\bm{a}(c)$ (nats),
$f(c)=-\min(\ell_\text{CTC}(c)/|\bm{y}_i|,10)\in[-10,0]$ is the negative CTC loss per character of the block's text
$\bm{y}_i$ ($|\bm{y}_i|$ characters) under an independent recogniser~\citep{baevski2020wav2vec}, and $\lambda\ge0$
converts unit-seconds of delay into the same log-score scale; the weights 0.5 and 5 were fixed before training and are
not tuned. A plan scores the sum of its potentials, and normalising the exponentiated score over all feasible plans gives
\begin{equation}
\begin{aligned}
p_\lambda(\bm{z}\mid\bm{x})&=\frac{\mathds{1}[\bm{z}\in\mathcal{Z}(\bm{x})]}{Z_\lambda(\bm{x})}\exp\sum_{i=1}^{I}\phi^\lambda_i(z_i),\\
Z_\lambda(\bm{x})&=\sum_{\bm{z}\in\mathcal{Z}(\bm{x})}\exp\sum_{i=1}^{I}\phi^\lambda_i(z_i),
\end{aligned}
\label{eq:crf}
\end{equation}
a distribution over plans with the scalar partition function $Z_\lambda(\bm{x})>0$. Writing feasibility as transition
factors, $\mathds{1}[\bm{z}\in\mathcal{Z}(\bm{x})]=\mathds{1}[h_1=0]\prod_{i<I}\mathds{1}[h_{i+1}=g_i(h_i,c_i)]$, shows
that Eq.~\eqref{eq:crf} is a linear-chain CRF~\citep{lafferty2001crf} whose node potentials are Eq.~\eqref{eq:plpot}
and whose transition potentials are $0$ or $-\infty$: the stream's first-in-first-out queue.

\paragraph{Exact inference.} $Z_\lambda$ sums over $\prod_i|C_i|=8^I$ candidate sequences. The chain structure makes
this tractable: the past influences the future only through the backlog. Group the partial sums by it, defining for
every block two non-negative vectors over the $|\mathcal{H}|$ backlog values,
$\alpha_i,\beta_i\in\R_{\ge0}^{|\mathcal{H}|}$:
\begin{equation}
\begin{aligned}
\alpha_i(h)&=\textstyle\sum_{c_{<i}:\,h_i=h}\exp\sum_{j<i}\phi_j(z_j),\\
\beta_i(h)&=\textstyle\sum_{c_{\ge i}\mid h_i=h}\exp\sum_{j\ge i}\phi_j(z_j),
\end{aligned}
\label{eq:msg}
\end{equation}
the total weight of feasible prefixes that leave block $i$ with backlog $h$, and of all completions from there.
Splitting off one block gives the recursions, with $\alpha_1(h)=\mathds{1}[h=0]$ and $\beta_{I+1}\equiv1$,
\begin{equation}
\begin{aligned}
\alpha_{i+1}(h')&=\textstyle\sum_{(h,c):\,g_i(h,c)=h'}\alpha_i(h)\,e^{\phi_i(h,c)},\\
\beta_i(h)&=\textstyle\sum_{c\in C_i}e^{\phi_i(h,c)}\,\beta_{i+1}\big(g_i(h,c)\big),
\end{aligned}
\label{eq:fb}
\end{equation}
so $Z_\lambda=\beta_1(0)$. A state $(h,c)$ at block $i$ is reached with weight $\alpha_i(h)$, contributes
$e^{\phi_i(h,c)}$ and is continued with weight $\beta_{i+1}(g_i(h,c))$; its posterior marginal, collected in a matrix
$\mu_i\in[0,1]^{|\mathcal{H}|\times|C_i|}$ (rows: backlog values, columns: candidates; entries sum to one), and the
expected delay of the stream are therefore
\begin{equation}
\begin{aligned}
\mu_i(h,c)&=\alpha_i(h)\,e^{\phi_i(h,c)}\,\beta_{i+1}\big(g_i(h,c)\big)/Z_\lambda,\\
\mathbb{E}_{p_\lambda}[\Phi]&=\textstyle\sum_{i,h,c}\mu_i(h,c)\,L_i(h,c),
\end{aligned}
\label{eq:marg}
\end{equation}
with $\Phi(\bm{z})=\sum_iL_i(z_i)$ the plan's total delay in unit-seconds. All of this is exact and costs
$O(I\,|\mathcal{H}|\max_i|C_i|)$ per stream (Proposition~\ref{prop:fb}, Appendix~\ref{app:proofs}); exactness holds
within the state space, whose candidates are finite samples and whose backlog is discretised.

\paragraph{Choosing $\lambda$: a latency budget.} Rather than tune $\lambda$, we ask for the distribution $q$ over the
plans of a training corpus $\mathcal{S}$ (a set of streams with $W$ target units in total) that is closest to the
quality-only CRF $p_0=p_{\lambda=0}$ while its expected delay per unit stays within a budget of $b$
seconds~\citep{ganchev2010posterior}; with $\Phi_\mathcal{S}=\sum_{\bm{x}\in\mathcal{S}}\Phi(\bm{z}_{\bm{x}})$ the
corpus delay in unit-seconds,
\begin{equation}
\min_q\ \mathrm{KL}(q\,\Vert\,p_0)\quad\text{s.t.}\quad\mathbb{E}_q[\Phi_\mathcal{S}]\le bW,
\label{eq:pr}
\end{equation}
Its Lagrangian is $\mathrm{KL}(q\Vert p_0)+\lambda(\mathbb{E}_q[\Phi_\mathcal{S}]-bW)$ with multiplier $\lambda\ge0$. For
fixed $\lambda$ the minimising $q$ is the Gibbs distribution $q_\lambda\propto p_0\,e^{-\lambda\Phi_\mathcal{S}}$.
Because $\Phi_\mathcal{S}$ is a sum of per-block delays, multiplying $p_0$ by $e^{-\lambda\Phi_\mathcal{S}}$ only
subtracts $\lambda L_i$ from every node potential, so $q_\lambda$ is exactly the CRF $p_\lambda$ of Eq.~\eqref{eq:crf},
one per stream with a shared $\lambda$; this is why the price of delay appears in Eq.~\eqref{eq:plpot}. Substituting
$q_\lambda$ back gives the scalar dual function
\begin{equation}
\mathcal{D}(\lambda)=-\sum_{\bm{x}\in\mathcal{S}}\log\frac{Z_\lambda(\bm{x})}{Z_0(\bm{x})}-\lambda bW,
\label{eq:dual}
\end{equation}
and since $\partial_\lambda\log Z_\lambda=-\mathbb{E}_{p_\lambda}[\Phi]$,
\begin{equation}
\begin{aligned}
\mathcal{D}'(\lambda)&=\mathbb{E}_{p_\lambda}[\Phi_\mathcal{S}]-bW,\\
\mathcal{D}''(\lambda)&=-\mathrm{Var}_{p_\lambda}[\Phi_\mathcal{S}]\le0.
\end{aligned}
\label{eq:dualderiv}
\end{equation}
The dual is concave and the expected delay falls monotonically in $\lambda$, so the optimum $\lambda^\star$ is $0$ when
$p_0$ already meets the budget and otherwise the root of $\mathbb{E}_{p_\lambda}[\Phi_\mathcal{S}]=bW$, which bisection
finds with Eq.~\eqref{eq:marg} (Proposition~\ref{prop:pr}). We set $b$ to the delay per unit of the base rate, without
the CRF, minus 0.3\,s, giving $\lambda^\star=0.50$ (Zh$\rightarrow$En) and $0.70$ (En$\rightarrow$Zh).

\paragraph{From the posterior to a Talker.} The CRF needs the release times of future blocks, which are unknown at
test time; we therefore distil its posterior into the Talker once, during training. On a training stream, condition the
marginal of Eq.~\eqref{eq:marg} on the backlog $\tilde{h}_i\in\mathcal{H}$ reached at block $i$ under the base rate; the
forward weight $\alpha_i(\tilde{h}_i)$ cancels and leaves a probability vector over the $|C_i|$ candidates,
$w_i\in[0,1]^{|C_i|}$,
\begin{equation}
\begin{aligned}
w_i(c)&=\frac{\mu_i(\tilde{h}_i,c)}{\sum_{c'}\mu_i(\tilde{h}_i,c')}\\
&=\frac{e^{\phi_i(\tilde{h}_i,c)}\,\beta_{i+1}\big(g_i(\tilde{h}_i,c)\big)}{\beta_i(\tilde{h}_i)},
\end{aligned}
\label{eq:w}
\end{equation}
which weighs a candidate by its own potential \emph{and}, through $\beta_{i+1}$, by the queueing cost it imposes on all
later blocks. The PA Talker $\pi_{\bm{\omega}}$ (parameters $\bm{\omega}$), which reads the block's text, the voice
prompt, the history and a backlog token $\langle\tilde{h}_i\rangle$ (one of 8 bins of 0.5\,s), is fine-tuned on the
weighted candidates:
\begin{equation}
\begin{aligned}
\mathcal{J}(\bm{\omega})&=\textstyle\sum_{\bm{x},i}\Big[-\sum_{c\in C_i}w_i(c)\log\pi_{\bm{\omega}}\big(\bm{a}(c)\mid\cdot\big)\\
&\quad+\kappa_\text{rate}\,\mathrm{H}\big(\rho_i,\pi_{\bm{\omega}}(\tau\mid\cdot)\big)\\
&\quad+\kappa_\text{kl}\,\mathrm{KL}\big(\pi_0\Vert\pi_{\bm{\omega}}\big)\Big],
\end{aligned}
\label{eq:distill}
\end{equation}
where $\pi_{\bm{\omega}}(\bm{a}(c)\mid\cdot)$ is the probability of $c$'s token sequence;
$\rho_i\in[0,1]^{35}$ is the posterior distribution over the Talker's 35 speaking-rate tokens,
$\rho_i(\tau)=\sum_cw_i(c)\mathds{1}[\tau(c)=\tau]$ with $\tau(c)$ the rate token of candidate $c$, and
$\mathrm{H}(\rho_i,\cdot)$ the cross-entropy of the Talker's rate prediction against it; $\pi_0$ is the base Talker (the KL keeps the codes close to
it); and $(\kappa_\text{rate},\kappa_\text{kl})=(1,0.5)$ are scalar weights. Writing
$\tilde{\pi}_{\bm{\omega}}(c)=\pi_{\bm{\omega}}(\bm{a}(c))/\pi_{\bm{\omega}}(C_i)$ for the Talker renormalised on the
candidates, the first term equals $\mathrm{KL}(w_i\Vert\tilde{\pi}_{\bm{\omega}})+\mathrm{H}(w_i)-\log\pi_{\bm{\omega}}(C_i)$
(Proposition~\ref{prop:distill}): it fits the Talker's preferences among the renderings to the CRF posterior, one
M-step of posterior-regularised EM. We run one round (Appendix~\ref{app:plem}).

\paragraph{Decoding.} At test time the PA Talker reads the backlog token of its actual queue, samples two candidates
per block and keeps the one with the higher $f$; the budget constrains the training posterior, not the online delay,
which we measure. The \emph{playback layer} is read-level merging (\S\ref{sec:method-online}), the backlog-dependent
base rate and the PA Talker; \sys{} is the deadline Thinker with this layer.

\subsection{Running on a whole meeting}
\label{sec:method-online}

\paragraph{Online windows.} A window opens at the first frame where any diarizer slot is active and closes after
0.6\,s with no active slot or at 20\,s; commits wait until the diarizer has released the frames that open (and, at the
last read, close) the window (\emph{chunk protocol}, Appendix~\ref{app:chunk}). Slots persist across windows; the
decoder restarts at each window. Blocks of one speaker committed in the same read are merged before they reach the
Talker (\emph{read-level merging}, rl).
